%! Inverse Functions — Unit 1, Lesson 1.6 — Guided Notes
\documentclass[10pt]{article}
\usepackage{precalculus-article}
\usepackage{precalculus-boxes}
\newcommand{\TallMath}[1]{$\displaystyle #1\rule[-1.4em]{0pt}{3.2em}$}

\begin{document}

\pageheader{Unit 1, Lesson 1.6}{Guided Notes}
\namedateperiod

\vspace{0.06in}

\begin{objectivebox}
By the end of this lesson, I will be able to\ldots
\begin{enumerate}[itemsep=2pt]
  \item Determine the \blank{4cm} of a function from a table or set of ordered pairs, and
        state when a function is \blank{3cm} on a given domain. \hfill\textit{(LO 2.8.A)}
  \item Find the inverse of a function \blank{3cm} by swapping $x$ and $y$ and solving,
        and verify the result using \blank{3cm}. \hfill\textit{(LO 2.8.B)}
  \item Restrict the \blank{3cm} of a non-invertible function to make it invertible.
        \hfill\textit{(LO 2.8.B)}
\end{enumerate}
\end{objectivebox}

\vspace{0.04in}

\begin{vocabbox}
\termblanklong{Inverse function}
\termblanklong{Invertible}
\termblanklong{One-to-one function}
\termblanklong{Domain restriction}
\termblanklong{$f^{-1}$ (f-inverse)}
\termblanklong{Identity function}
\termblanklong{Reflection over $y = x$}
\end{vocabbox}

\vspace{0.04in}

\begin{hookbox}
A nurse uses the formula $D(w) = 1.5w + 10$ to find a patient's medication dose $D$ (mg)
from their weight $w$ (kg).

\medskip
If a patient received 55 mg, what was their weight? \blank{5cm}

\medskip
The formula that converts dose back to weight is called the \blank{5cm} of $D$.
We write it as \blank{3cm}.
\end{hookbox}

\vspace{0.04in}

\begin{notesbox}{LO 2.8.A --- Inverse Functions from Tables (EK 2.8.A.1 \& 2.8.A.2)}
\textbf{Key Idea:} If $f(a) = b$, then $f^{-1}(b) = $ \blank{2cm}.
The inverse reverses every \blank{4cm} pair.

\medskip
\textbf{Step 1: Complete the output table for $f(x) = 4x - 1$.}

\renewcommand{\arraystretch}{1.7}
\begin{tabular}{|c|c|c|c|c|}
\hline
$x$    & 1 & 2 & 3 & 4 \\
\hline
$f(x)$ & \blank{1.5cm} & \blank{1.5cm} & \blank{1.5cm} & \blank{1.5cm} \\
\hline
\end{tabular}

\medskip
\textbf{Step 2: Fill in the inverse table by reversing every pair.}

\renewcommand{\arraystretch}{1.7}
\begin{tabular}{|c|c|c|c|c|}
\hline
$x$          & \blank{1.5cm} & \blank{1.5cm} & \blank{1.5cm} & \blank{1.5cm} \\
\hline
$f^{-1}(x)$  & \blank{1.5cm} & \blank{1.5cm} & \blank{1.5cm} & \blank{1.5cm} \\
\hline
\end{tabular}

\medskip
\textbf{Invertible check:} Does $f$ map each input to a \blank{4cm} output?
If any two inputs share the same output, the function is \textbf{not} invertible. \blank{3cm}

\medskip
Domain and range swap: the domain of $f^{-1}$ is the \blank{3cm} of $f$, and vice versa.
\end{notesbox}

\vspace{0.04in}

\begin{notesbox}{LO 2.8.B --- Analytic Inverse Method (EK 2.8.B.4)}
To find $f^{-1}(x)$ for $f(x) = 4x - 1$:

\medskip
\textbf{Step 1:} Write $y = 4x - 1$.

\medskip
\textbf{Step 2:} Swap $x$ and $y$: \quad $x = $ \blank{5cm}

\medskip
\textbf{Step 3:} Solve for $y$:
\[
  x + 1 = 4y \qquad \Rightarrow \qquad y = \dfrac{x + 1}{4}
\]

\medskip
\textbf{Step 4:} Write $f^{-1}(x) = $ \blank{4cm}
\end{notesbox}

\vspace{0.04in}

\begin{notesbox}{Composition Verification (EK 2.8.B.1)}
\textbf{Key property:} $f(f^{-1}(x)) = x$ and $f^{-1}(f(x)) = x$.

\medskip
\textbf{Verify $f(f^{-1}(x)) = x$ for $f(x) = 4x - 1$ and $f^{-1}(x) = \dfrac{x+1}{4}$:}

\[
  f\!\left(f^{-1}(x)\right) = 4 \cdot \left(\rule{3cm}{0.4pt}\right) - 1
  = \rule{4cm}{0.4pt} = x \quad \checkmark
\]

\medskip
\textbf{Also verify $f^{-1}(f(x)) = x$:}

\[
  f^{-1}\!\left(f(x)\right) = \dfrac{(\rule{3cm}{0.4pt}) + 1}{4}
  = \rule{4cm}{0.4pt} = x \quad \checkmark
\]

\medskip
Both compositions must equal $x$; if either fails, the inverse is \blank{3cm}.
\end{notesbox}

\vspace{0.04in}

\begin{notesbox}{Domain Restriction (EK 2.8.A.1, EK 2.8.B.5)}
$f(x) = x^2$ is \textbf{not} invertible on all real numbers because:
$f(2) = $ \blank{2cm} and $f(-2) = $ \blank{2cm} --- two different inputs produce the same output.

\medskip
Restrict the domain to $x \ge 0$: then every output comes from exactly \blank{2cm} input.

\medskip
On the restricted domain $[0, \infty)$: $f^{-1}(x) = $ \blank{4cm}

\medskip
The domain of $f^{-1}$ is \blank{3cm}; the range is \blank{3cm}.
\end{notesbox}

\end{document}
